\section{A numerical detector bound for Gaussian determinant three}
\label{sec:e8-detector-degree}

\paragraph{Check status.}
The entire proof, including the incidence-component and rational
image-span arguments, has been read by the author and two independent
readers. Their actual events are recorded in CHECK\_EVENTS.tsv.
The geometric, Fermat and lattice inputs were checked separately.
D-E8-XFER01--03 proves the tensor-diagram lemma.

\subsection*{HD-E8Q1. Constants and the bound}

All degrees below are intersection degrees for the stated integral
polarizations. Put
\[
\begin{gathered}
 G=1393459200,\quad g=1+53G/24=3077222401,\quad M=2G,\\
 R=(2G+1)G!G^G,\qquad L=2gR,\\
 A=(4L)!\,[1000g(G+1)(L+1)]^{4L},\qquad
 B=(4L+4)A^{2L+4},\\
 S=[26(L+1)]^L\mathcal H_{12}(L),\qquad
 D_{\rm src}=M^6B^{4L+21}S.
\end{gathered}                                                    \tag{E8Q.1}
\]
Here the explicit integer \(\mathcal H_{12}\) is defined in
SF-SKQ01--05. Every constant uses only factorials, products and
powers of positive integers. The constants are deliberately generous.

For a maximal-order polarized Gaussian sixfold of signature
\((3,3)\), signed determinant class three, and polarization volume
\(\Delta\), define
\[
 \boxed{D_3(\Delta)=69984^{12}M^{114}\Delta^{18}B^{4L+21}S.}       \tag{E8Q.2}
\]
For the fixed order \(\mathbb Z+f\mathbb Z[i]\), replace
\(\Delta\) by \(f^{12}\Delta\), giving a factor \(f^{216}\).
The assertions to prove are
\[
 \delta_W\le D_3(f^{12}\Delta),\qquad
 c_t\le5462\,\delta\lceil\sqrt T\rceil
              (5f^2)^{24}D_3(f^{12}\Delta)^2.                   \tag{E8Q.3}
\]
The normalization of \(\delta,T,c_t\) is fixed in the introduction.
The second estimate uses HD-U with \(a=f(2+i)\), after the first
has supplied a cycle.

\subsection*{HD-E8Q2. The full Brauer identity and its slots}

We need an identity of representations, rather than only an identity
of determinant characters. D-E8-XFER01 applies the integer
column-lattice argument from HECKE\_DIV SQ3 to all irreducible
representation multiplicities. It gives
\[
 \rho\oplus R_-\simeq R_+,\qquad
 R_\pm=\bigoplus_j
       (\operatorname{Ind}_{H_j}^G\chi_j)^{n_{j,\pm}},\qquad
 \sum_j(n_{j,+}+n_{j,-})\le R.                                 \tag{E8Q.4}
\]
The \(\chi_j\) include trivial characters when necessary. This
uses Brauer's integral induction theorem and the proved integer
minor/Cayley-graph bound, not an inference from determinant
characters to a representation identity.

Let \(V_j\) be the associated character multiplicity in \(H^1(C)\),
of dimension \(v_j\le2g\), repeated according to its coefficient.
Replace every minus factor by its conjugate character space.
The number of degree-one curve slots is
\[
 l=\sum_jv_j\le2gR=L,\qquad
 l=6+2\dim\operatorname{Hom}_G(R_-,H^1(C)).                     \tag{E8Q.5}
\]
Thus \(l\) is even and at least six. Zero-dimensional character
spaces have unit determinant and are omitted. The determinant
tensor of these suppliers contracts nontrivially to the primitive
sixfold determinant by (E8Q.4) and polarization duality.
D-E8-XFER02--03 realizes this contraction by bounded algebraic
tensor diagrams. It never takes a root of a determinant cycle.

\subsection*{HD-E8Q3. Degrees of the returning cyclic correspondences}

Choose a line bundle \(H_C\) of degree \(6g\) on \(C\).
It is very ample, since \(6g\ge2g+1\). Give every curve product
the sum polarization; its numerical class is invariant under
the deck automorphisms.

For a nontrivial character in (E8Q.4), let
\(C/\ker\chi_j\to B_j=C/H_j\) be its faithful cyclic cover.
Put \(b_j=g(B_j)\), let \(r_j\) be its number of branch points,
and write \(v_j=2b_j-2+r_j\). Every intermediate inertia order
divides twelve. When \(r_j\ge3\), HD-FS1 uses
\[
 Y_j=\mathbf P_{F_{12}^{r_j-2}}
       (\mathcal O^{b_j}\oplus\mathcal O(-12)),\qquad
 d_j=\dim Y_j=b_j+r_j-2\le v_j.                               \tag{E8Q.6}
\]
Use the convention of lines, write
\(\xi_j=c_1(\mathcal O_{Y_j}(1))\), let \(x_j\) be the
Fermat hyperplane class, and put \(H_j=\xi_j+x_j\).
This is very ample: its section bundle on the base is
\(\mathcal O(1)^{b_j}\oplus\mathcal O(13)\), a sum of
very ample line bundles. Their sections separate base points
and base tangent vectors using any nonzero fibre coordinate;
the full linear system on each projective fibre separates its
points and tangent vectors. This proves the required embedding.
Both \(\xi_j\) and \(x_j\) are globally generated and nef.

For \(r_j=2\), take one component \(Y_j=\mathbf P^{b_j}\).
For \(r_j=0\), take \(Y_j=\mathbf P^{b_j-1}\). Use their
hyperplane classes. These are the split projective suppliers
of HD-FS1, and \(d_j\le v_j\) still holds. The case \(v_j=0\)
was discarded. A zero-dimensional Fermat scheme is treated
componentwise, so its number of points does not multiply an
individual degree bound.

Put \(L_j=K_{B_j}(D_j)\) in the branched case and
\(L_j=K_{B_j}\) in the unramified case. Then
\(\deg L_j=v_j\), and the morphism \(Y_j\to|L_j|\)
pulls back the hyperplane class to \(\xi_j\), or to the
specified projective-space hyperplane class.

In \(C^{v_j}\times Y_j\), require the section represented by
\(y\in Y_j\) to vanish at the image of each curve coordinate
in \(B_j\). There are \(v_j\) Cartier conditions, with classes
\[
 \operatorname{pr}_{C,a}^*q_j^*c_1(L_j)
                     +\operatorname{pr}_{Y_j}^*\xi_j,          \tag{E8Q.7}
\]
where \(q_j:C\to B_j\) has degree at most \(G\).
Thus \(\deg q_j^*L_j\le Gv_j\le2Gg\). Every class in
(E8Q.7) is bounded numerically by \(G+1\) times the total
product polarization, with nef difference.

This incidence is finite flat over \(Y_j\). The nonzero
universal section on the smooth curve has a zero divisor of
constant degree \(v_j\); its pullback to \(C\) also has
constant degree. The relative Cartier divisor is therefore
finite flat, and the independent coordinate fibre products
preserve this property. Repeated zeros produce finite
multiplicities, without excess-dimensional components.

The returning correspondence in HD-FS1 is a component of
this incidence: its ordered images form exactly the divisor
of the chosen section, with the chosen residual cyclic label.
A component assertion can be checked over a dense open set.
The system \(|L_j|\) is basepoint-free: in the branched
case its degree is at least \(2b_j\), and in the unramified
case it is the canonical system for \(b_j\ge2\).
Its general divisor is reduced and avoids the finitely many
branch points of \(C\to B_j\). The supplier map to
\(|L_j|\) is dominant, so these conditions define a dense
open in \(Y_j\). There the incidence is finite etale.
Requiring all projected zeros to be distinct and the
residual label to equal the chosen lift selects an
open-and-closed union of its components. Choosing one
component and taking its closure gives the desired
irreducible returning correspondence on the entire source.
A resolution pushes to its fundamental cycle, so its degree
is bounded by that of the whole incidence. This controls the
actual returning cycle without assigning a degree to an
unspecified resolution.

The bundle relation gives
\[
 \pi_*\xi_j^{b_j+k}=12^kx_j^k\quad(k\ge0),\qquad
 \deg_{H_j}Y_j
  =12\sum_{a=0}^{r_j-2}\binom{d_j}{a}12^{r_j-2-a}
  \le12\,13^{d_j}.                                            \tag{E8Q.8}
\]
For projective spaces the degree is one. A product of these
spaces of total dimension \(d\le L\) consequently has degree
at most \(d!\,12^L13^L\). A product with at most \(L\)
curve slots has dimension at most \(2L\) and degree at most
\((2L)!(6g)^L12^L13^L\). Intersecting the incidence equations
multiplies this upper bound by at most \((G+1)^L\).

For a trivial character, let \(b=g(C/H)\). Its determinant
has \(v=2b\) slots and is Tate. A product of \(b\) diagonals
on \((C/H)^{2b}\), followed by alternation and the
degree-one Kunneth projector, has nonzero top exterior
component: it is the \(b\)-th exterior power of the
nondegenerate symplectic tensor. Pull it to \(C^{2b}\).
Each factor is \(C\times_{C/H}C\), an effective cycle
of degree at most \(12gG\). The product degree is at most
\(b!(12gG)^b\). It is an external Tate cycle, not a
finite correspondence over a point.

Take the product of the nontrivial returning correspondences
and these Tate cycles. If the trivial factors occupy \(l_0\)
slots, the correspondence from \(Y=\prod Y_j\) to \(C^l\)
has dimension \(d+l_0/2\). The input cohomological degree
is \(2d-(l-l_0)\), as in HD-FS1; it is even because
\(l\) and \(l_0\) are even. All factorials from products
and incidence degrees are bounded by
\((2L)!^2\le(4L)!\); all the displayed degree powers are
bounded by \([1000g(G+1)(L+1)]^{4L}\). Hence each returning
correspondence has degree at most \(A\). This bound also
dominates every ambient product degree used here.

\subsection*{HD-E8Q4. Effective Hodge generators on the bundle suppliers}

Let \(F\) be the product of the positive-dimensional Fermat
bases of \(Y\), of total dimension \(t\le d\le L\).
Projective-space factors are bundles over a point. The
projective bundle formula decomposes the full cohomology as
\[
 \xi_1^{j_1}\cdots\xi_k^{j_k}\pi^*H^*(F),
                    \qquad 0\le j_i\le b_i,                  \tag{E8Q.9}
\]
where \(b_i\) is the projective fibre dimension. These
algebraic maps respect rational Hodge structures.
SF-SKQ gives finite effective integral spanning cycles
\(W\) on \(F\), each of degree at most
\(\mathcal H_{12}(t)\le\mathcal H_{12}(L)\).
For a point use its degree-one class.

Flat pullback of \(W\), intersected with general members
of the globally generated \(\xi_i\) systems to the powers
in (E8Q.9), is an effective integral cycle, possibly zero,
with the required class. Give \(Y\) its sum polarization
\(H_Y=\sum_i(\xi_i+x_i)\), omitting \(x_i\) for projective
spaces. To bound the degree, expand the required power of
\(H_Y\) and push to \(F\). A surviving term uses
\(\pi_*\xi_i^{b_i+a}=12^a x_i^a\); for projective spaces
only \(a=0\) survives. Since \(j_i\le b_i\), the factor
\(12^a\) is bounded by assigning a factor twelve to each
\(\xi_i\) selected from this expansion. The sum of all
nonnegative coefficients is therefore at most
\((13k)^d\), with \(k\le L\), or one for a point.
Each remaining monomial in the nef \(x_i\), of total
degree \(\dim W\), has intersection with \(W\) at most
\(\deg_{\sum x_i}W\), by expansion of that total power.
Consequently all these generators have degree at most
\[
                   [26(L+1)]^L\mathcal H_{12}(L)=S.           \tag{E8Q.10}
\]
They span every rational Hodge space of \(Y\), in all
codimensions, by (E8Q.9).

The Brauer supplier tensor has type \((l/2,l/2)\):
the sixfold determinant has type \((3,3)\), and the
minus determinant together with its conjugate contributes
a Tate factor. To justify rational descent, use a finite
Galois splitting field containing all \(G\)-th roots of
unity and the exact primitive root-line projector as in
D-E8-XFER02. Expand all coefficient-field projectors
and transfer maps into their integral correspondence terms.
Let \(I\) be the sum of the rational cohomological images
of those terms, before or after the fixed final map to the
sixfold. These are maps of pure Hodge structures with the
required Tate shifts, so \(I\) is a rational Hodge substructure.
The extended image contains the nonzero determinant line
for the natural character; Galois invariance then includes
its conjugate. Their rational span is the Weil plane.
In particular \(I\) contains a nonzero rational Weil class.
Semisimplicity of polarizable pure Hodge structures gives
a rational Hodge preimage in the appropriate source
cohomology. The effective generators (E8Q.10) span that
entire Hodge space, so some combination has nonzero final
Weil projection. No individual unbalanced character
determinant is asserted to be a Hodge class.

\subsection*{HD-E8Q5. Degree costs and an integral detecting component}

Use the published moving estimate
\cite[Lemma~2.2, author-version p.~13]{Truong2020}.
On a smooth projective ambient \(V\) of dimension \(m\),
the safe factor \(2(m+1)\deg(V)^{m+1}\) represents an
intersection of two effective integral cycles by a
difference of effective integral cycles, each with degree
at most that factor times the two input degrees.

Suppose a correspondence \(\Gamma\subset X\times X'\)
and its smooth product ambient have degrees at most \(A\),
and the ambient dimension is at most \(2L\). If an effective
cycle \(W\subset X\) has degree \(E\), then
\(\deg(W\times X')\le AE\). Indeed the Segre binomial
factor for \(W\times X'\) is no larger than that for
\(X\times X'\), and \(\deg X\ge1\).
The moving estimate, intersection with \(\Gamma\), and
proper projection express \(\Gamma_*[W]\) as a difference
of effective integral cycles with individual degrees at most
\[
                  2(2L+1)A^{2L+3}E\le BE.                   \tag{E8Q.11}
\]
Projection does not increase degree: the target hyperplane
pullback is nef and bounded by the total product hyperplane.
Components whose images drop dimension push to zero.
This estimate allows improper intersections at the outset.

All required correspondences meet these hypotheses.
HD-E8Q3 bounds the returning operation. The degree-one
Kunneth projector on \(C\) is
\(\Delta_C-[c]\times C-C\times[c]\); each term has degree
at most \(12g\). Its \(l\)-fold tensor has terms of degree
at most \(l!(12g)^l<A\).
By D-E8-XFER02 the determinant transfer is a sum of products
of \((l-6)/2+6\) deck graphs in \(C^{l+6}\), each of degree
\[
           ((l-6)/2+6)!(12g)^{(l-6)/2+6}<A.                  \tag{E8Q.12}
\]
Its cyclic-character projectors, representation matrix units,
and determinant permutations are already expanded into these
diagrams. The last primitive projector similarly expands
into coordinatewise deck graphs. All curve-product ambients
have dimension at most \(2L\) and degree at most \(A\).

There are at most four grouped correspondence operations:
the returning map; the full tensor Kunneth projector; the
transfer diagram; and the final primitive projector.
Thus (E8Q.11) bounds individual effective generators at
the end by \(B^4S\le B^{4L+21}S\).
The large exponent in (E8Q.1) is only a convenient
conservative bound; no uncounted tensor operation is needed.

Here is the precise coefficient and effectivity argument.
At every moving step retain the two effective parts as
separate generators. Their rational spans contain the
image of the previous span. Expand every coefficient-field
linear combination into its individual algebraic diagrams.
The final span has nonzero Weil projection after the fixed
last map to the sixfold. If all its integral effective
generators had zero such projection, so would the entire
span. Thus one generator, and then one irreducible component,
has nonzero projection. Taking that component does not
increase degree. Intermediate individual terms need not
belong to the intermediate projector images.
No integral saturation or division of an effective cycle
by a rational coefficient is claimed.

\subsection*{HD-E8Q6. The sixfold map and the target bound}

Let \(p\) be the self-adjoint primitive projector of
D-E8-LAT02. Its image \(B_C=pJ(C)\) is the abelian sixfold,
polarized by \(\theta_C=\Theta_{J(C)}|_{B_C}\).
The map \(f=Mp:J(C)\to B_C\) is an actual homomorphism and
\[
                    f^*\theta_C\le M^2\Theta_{J(C)}          \tag{E8Q.13}
\]
as semipositive forms, since \(p\) is the orthogonal
projection for the Rosati metric. Compose six copies of
\(f\) with Abel maps and addition to obtain
\(a_6:C^6\to B_C\). Cauchy's inequality for addition
of six vectors gives
\[
 a_6^*\theta_C\le6M^2\sum_{j=1}^6a_j^*\Theta_{J(C)}
                     =M^2\sum_{j=1}^6H_{C,j}                \tag{E8Q.14}
\]
numerically: the Abel theta degree is \(g\), whereas
\(\deg H_C=6g\). Positive representatives for this
numerical equality are furnished by the pulled-back
theta forms.

On the tensor product of six degree-one factors, addition
induces exterior multiplication, nonzero on the primitive
determinant. Equivalently, pullback of the conjugate Weil
determinant has its nonzero alternating tensor as the
component with one degree-one class on each curve.
The injective map \(f^*:H^1(B_C)\to H^1(C)\) has
nondegenerate restricted pairing. Its determinant pairing
with the output of D-E8-XFER02 is nonzero. Projection
formula therefore proves nonzero Weil projection after
\(a_{6*}\), as used in the component selection above.
By (E8Q.14), a surviving three-cycle of curve-product
degree at most \(B^{4L+21}S\) pushes to degree at most
\[
                    M^6B^{4L+21}S=D_{\rm src}.              \tag{E8Q.15}
\]

This bound is independent of branch points.
D-E8-GRR and D-E8R each provide a connected component
with period rank nine, so the natural sixfolds cover an
analytic open subset of their fixed-data period domain.
Their polarization volume is at most \(M^6\), by the
actual integral projector calculation in D-E8-LAT02.

For a specified maximal-order target lattice of signed
determinant class three, D-E8-LAT01--02 supplies one fixed
rational Hermitian isometry and a clearing integer
\[
                      N\le69984^2M^{18}\Delta^3.
\]
On the identified marked period domains its isogeny
pulls the target polarization back to \(N^2\theta_C\).
Pushing a detector multiplies its degree by \(N^6\).
Thus an analytic open subset of the target family has
detectors bounded by
\[
 N^6D_{\rm src}
       \le69984^{12}M^{114}\Delta^{18}B^{4L+21}S=D_3(\Delta).
\]
The accepted bounded-cycle closedness and fixed-data
spreading input extends this integer bound to the entire
connected family. Source lattice data are locally constant
on the chosen Hurwitz component, and the same estimate
works for every one of them.

For a fixed conductor \(f\), use
\(f\mathbb Z[i]\Lambda\subset\Lambda\) as in D-E8-LAT02.
Its pulled-back polarization volume is at most
\(f^{12}\Delta\); the last identity isogeny pulls the
target polarization back exactly. This proves the first
inequality in (E8Q.3). Applying HD-U with \(n=3\) and
\(a=f(2+i)\), of norm \(5f^2\), proves the second.
The resulting division bound has conductor exponent
\(432+48=480\) and polarization-volume exponent \(36\).
Cycles are supplied before positivity and the integral
polynomial projector of HD-U are applied.
